\documentclass{article}

\usepackage{amsmath,amssymb}
\usepackage{xurl}
\usepackage[hidelinks]{hyperref}

% Shared commands for notes in docs/paper-gaps/.

\DeclareUnicodeCharacter{2080}{\ifmmode{}_{0}\else\textsubscript{0}\fi}
\DeclareUnicodeCharacter{2081}{\ifmmode{}_{1}\else\textsubscript{1}\fi}
\DeclareUnicodeCharacter{2082}{\ifmmode{}_{2}\else\textsubscript{2}\fi}
\DeclareUnicodeCharacter{2099}{\ifmmode{}_{n}\else\textsubscript{n}\fi}

\newcommand{\Lean}{\textsc{Lean}}
\ExplSyntaxOn
\NewDocumentCommand{\leanid}{m}{
  \ifmmode
    \textsf{\detokenize{#1}}
  \else
    \group_begin:
    \urlstyle{sf}
    \tl_set:Nn \l_tmpa_tl {#1}
    \tl_replace_all:Nnn \l_tmpa_tl {\_} {_}
    \exp_args:NV \path \l_tmpa_tl
    \group_end:
  \fi
}
\ExplSyntaxOff
\providecommand{\tr}{\operatorname{tr}}
\newcommand{\tnleanIssueBase}{https://github.com/LionSR/TNLean/issues/}
\newcommand{\tnleanPrBase}{https://github.com/LionSR/TNLean/pull/}
\newcommand{\tnleanDocBase}{https://sirui-lu.com/TNLean/docs/}
\newcommand{\ghissue}[1]{\href{\tnleanIssueBase#1}{GitHub issue \##1}}
\newcommand{\ghpr}[1]{\href{\tnleanPrBase#1}{PR \##1}}
\newcommand{\doclink}[1]{\href{\tnleanDocBase#1.html}{\path{#1}}}

% Dirac notation and the identity operator, as in blueprint/src/macros/common.tex.
% \providecommand keeps note-local definitions authoritative.
\usepackage{dsfont}
\providecommand{\ket}[1]{|#1\rangle}
\providecommand{\bra}[1]{\langle #1|}
\providecommand{\braket}[2]{\langle #1 | #2 \rangle}
\providecommand{\ketbra}[2]{|#1\rangle\!\langle #2|}
\providecommand{\Id}{\mathds{1}}
\providecommand{\one}{\mathds{1}}
\providecommand{\C}{\mathbb{C}}
\providecommand{\MN}[1]{M_{#1}(\C)}
\providecommand{\MD}{M_D(\C)}

% External referencing for paper sources.  The build helper writes sanitized
% auxiliary files under build/paper-aux/ and excludes duplicated source labels.
\usepackage{xr}

% Import a generated paper auxiliary file with a caller-supplied label prefix.
% The candidate paths support builds from docs/paper-gaps/, the repository root,
% and build/paper-aux/ itself.
\newcommand{\importpaperaux}[2]{%
  \IfFileExists{../../build/paper-aux/#2.aux}{%
    \externaldocument[#1]{../../build/paper-aux/#2}%
  }{%
    \IfFileExists{build/paper-aux/#2.aux}{%
      \externaldocument[#1]{build/paper-aux/#2}%
    }{%
      \IfFileExists{#2.aux}{%
        \externaldocument[#1]{#2}%
      }{}%
    }%
  }%
}

% General external reference macros
\newcommand{\paperref}[2]{\ref{#1-#2}}
\newcommand{\papereqref}[2]{(\ref{#1-#2})}

% CPSV16 source (1606.00608 / MPDO-22-12-17-2.tex) external document and shortcuts
\importpaperaux{cpsv16-}{1606.00608/MPDO-22-12-17-2-tnlean}
\newcommand{\cpsvref}[1]{\paperref{cpsv16}{#1}}
\newcommand{\cpsveqref}[1]{\papereqref{cpsv16}{#1}}

% Verdict marker: \gapnote{<kind>}{<status>}. Kind is the policy
% classification (clarification, local-correction, scope-restriction,
% unfaithful, false-source, open-gap); status is open, wip (elimination
% underway), resolved, or historical. Machine-read by the site index;
% prints a verdict line.
\newcommand{\gapnote}[2]{\par\noindent\textbf{Verdict:} #1 (#2).\par}


% Fallbacks keep this author document compilable when it is built outside its
% source directory and a different command.tex is found first.
\providecommand{\leanid}[1]{\textsf{\detokenize{#1}}}
\providecommand{\doclink}[1]{%
    \href{https://sirui-lu.com/TNLean/docs/#1.html}{\path{#1}}}
\providecommand{\gapnote}[2]{%
    \par\noindent\textbf{Verdict:} #1 (#2).\par}

\title{The Unbroken Subgroup Is the Stabilizer of a Block}
\author{}
\date{2026-09-27}

\begin{document}
\maketitle
\gapnote{local-correction}{resolved}

Let a finite group $G$ act by matrix product unitaries with three-cocycle
$\omega$ on the set $\mathcal I$ of blocks $x$ of a symmetry-broken ground space,
and assume, as both sources do, that the action on $\mathcal I$ is transitive.
The two sources define the unbroken subgroup $H$ differently, and only one of
the definitions supports the classification of phases by pairs $(H,\alpha)$
with $\alpha\in\mathcal H^2(H,U(1))$.

\section{The two definitions}

Garre-Rubio and Schuch define $H$ as the set of $h\in G$ with $h\cdot x=x$ for
every $x\in\mathcal I$, and state that the ground-state degeneracy is
$|G/H|$~\cite{GarreRubioSchuch2024DomainWall}.\footnote{Source traceability:
\path{Papers/2405.00439/MPU-DW.tex}, line 1880 and lines 2013--2036.} This
is the kernel
\begin{align}
    K=\bigcap_{x\in\mathcal I}\operatorname{Stab}(x)
    \label{eq:gs24_kernel}
\end{align}
of the action on blocks. Garre-Rubio, Lootens and Moln\'ar, whose
classification the first paper reviews, define $H$ as the subgroup of elements
that do not permute one chosen block $x_0$, that is, the stabilizer
$\operatorname{Stab}(x_0)$, and state that the number of blocks is
$|G/H|$~\cite{GarreRubio2023MPOSym}.\footnote{Source traceability:
\path{Papers/2203.12563/REsubmission.tex}, lines 734--761; the stabilizer is
defined at line 736, the two-cocycle $\psi_{x_0}$ at lines 748--749, and the
reconstruction of the L-symbols from $(H,\psi)$ at lines 755--760.}

For a transitive action $\operatorname{Stab}(g\cdot x_0)=g\operatorname{Stab}(x_0)g^{-1}$,
so $K$ is the normal core of $\operatorname{Stab}(x_0)$, the largest normal
subgroup of $G$ contained in it. The two definitions agree exactly when
$\operatorname{Stab}(x_0)$ is normal in $G$. This holds for every abelian $G$,
and therefore for every example printed in either source.

\section{Which definition the classification needs}

The classification needs the stabilizer, for three reasons.

\begin{enumerate}
    \item \emph{Counting blocks.} By the orbit--stabilizer theorem
          $|\mathcal I|=|G/\operatorname{Stab}(x_0)|$, while
          $|G/K|\geq|\mathcal I|$ with equality only when the stabilizer is
          normal. The degeneracy $|G/H|$ asserted at line 1880 of the first
          source holds for the stabilizer.
    \item \emph{Trivializing the cocycle.} The coupled pentagon equation
          $L^x_{h,k}L^x_{g,hk}=\omega(g,h,k)L^{k\cdot x}_{g,h}L^x_{gh,k}$,
          restricted to $g,h,k$ in a subgroup fixing the single block $x$,
          is the statement that $\omega$ restricted to that subgroup is a
          coboundary with potential $L^x$. Fixing one block suffices, so the
          argument gives the stronger conclusion for the larger subgroup
          $\operatorname{Stab}(x_0)\supseteq K$.
    \item \emph{Reconstruction.} The two-cocycle $\psi_{x_0}$ is defined on
          the stabilizer, and the L-symbols are rebuilt from $(H,\psi)$
          through a choice of coset representatives $k_y$ with
          $y=k_y\cdot x_0$ and the unique factorization $gk_x=k_yh$ with
          $h\in H$. This factorization is the coset decomposition of $G$ by
          $\operatorname{Stab}(x_0)$; with $K$ in its place the
          representatives do not correspond to blocks. The pairs $(H,\psi)$,
          with $H$ determined up to conjugacy, are the data classifying the
          indecomposable module categories of $\mathrm{Vec}_G^\omega$.
\end{enumerate}

The domain-wall discussion of the first source also uses the stabilizer
reading for normal $H$: when $\operatorname{Stab}(x_0)$ is normal, every
block has the same stabilizer, so $g\notin H$ implies $g\cdot x\neq x$ for
every block $x$, as assumed at line 2013 of that source.

\section{A non-normal example}

Let $G=S_3$ act on $\mathcal I=\{1,2,3\}$ by permutations, realized by the
on-site permutation representation of $S_3$ on $\mathbb C^3$ acting on the
three product states $|x\rangle^{\otimes N}$; the three-cocycle is trivial.
The action is transitive, and
\begin{align}
    \operatorname{Stab}(3) & =\{e,(12)\}, &
    K                      & =\{e\}.
    \label{eq:gs24_s3_subgroups}
\end{align}
The stabilizer has order two and is not normal, since
$(13)(12)(13)=(23)\notin\operatorname{Stab}(3)$. The stabilizer reading
gives three blocks, $|S_3/\operatorname{Stab}(3)|=3$, which is correct; the
kernel reading gives $|S_3/K|=6$. The transposition $(12)\notin K$ fixes the
block $3$, so a group element outside the kernel may still act on some block
as an unbroken symmetry.

\section{The formal statement}

The formal theorem that a subgroup fixing a block trivializes the
restricted three-cocycle takes an arbitrary homomorphism $f:H\to G$ and one
block $x$ with $f(a)\cdot x=x$ for all $a\in H$, that is, the stabilizer
reading.\footnote{Formal declaration:
\leanid{MPOTensor.GroupFamily.IsNormalRepresentation.isTrivialGaugeClass_comap_omega_of_fixed}
in \doclink{TNLean/MPS/Symmetry/MPOSymmetry/PermutedBlocks}.} It applies in
particular to the kernel, which fixes every block.

\textbf{Verdict: resolved local correction.} In the classification of
arXiv:2405.00439, the unbroken subgroup $H$ is read as the stabilizer of one
block, as in arXiv:2203.12563; the printed definition as the kernel of the
action agrees with it whenever the stabilizer is normal, in particular for
every abelian group.

\bibliographystyle{alpha}
\bibliography{references}

\end{document}
